% ============================================================================
\clearpage
\section*{v4.10.0-r59 Source-Only: The Handedness Series and the Hilbert-Shift Tower (k245--k274, H125, H126)}
\addcontentsline{toc}{section}{v4.10.0-r59 Source-Only: The Handedness Series and the Hilbert-Shift Tower (k245--k274, H125, H126)}
\label{sec:v4100-r59-handedness-tower}
% ============================================================================

\paragraph{Purpose.}
Between 23 and 24 November 2026 a series of explorations and two tests followed r58: where the handedness (the sign \(\tau\varepsilon\) of
r58.4) can come from, and whether the field value of the photon-clock picture has a generation rule from one Hilbert shift to the next.
They were recorded as research notes with scripts and outputs, not as layers.  The author's standing rule that a finding used in an article
enters the Companion as a layer (``the main paper stays the source'') requires them here before the German preprint cites them.  This layer
therefore collects them, each with its status as recorded at the time.  It adds no new computation; every statement below is checked in the
cited scripts.  All items are explorations and \emph{not blind}, except H125 and H126, which were fixed in advance notes before the computation.
Earlier layers are not rewritten; where an item rereads an earlier layer (r13, r14.4, r43) the earlier text stands.

\noindent\textbf{Source-only status.}
This v4.10.0-r59 layer is source-only.  It does not build a PDF, does not assign a new DOI,
and does not authorize a public release.  The v5.3.0 Version DOI
\href{https://doi.org/10.5281/zenodo.24434403}{10.5281/zenodo.24434403} and the Concept DOI
\href{https://doi.org/10.5281/zenodo.19210728}{10.5281/zenodo.19210728} are unchanged by it.

\subsection*{v4.10.0-r59 setting}
\label{subsec:v4100-r59-setting}

Part A (r59.1--r59.6) follows the handedness: the orientation of the octant step, which r31 identified as a time orientation of the readout
and r58 as the product \(\tau\varepsilon\) of the time direction \(\tau\) of the readout and the orbit orientation \(\varepsilon\).
Part B (r59.7--r59.14) follows the Hilbert-shift tower HS\((n)\): the field value \(\psi\) lives in the Cayley--Dickson algebra \(A_n\) of
dimension \(2^n\) (the carrier sizes of the atlas), each shift adjoins one imaginary unit, and the basic orbit is
\(\psi(t)=\exp(u\omega t)\psi_0\) with a unit imaginary \(u\).  The assignment of the algebras to the shifts is a reading; the algebra itself is
standard and only recomputed.  Frame convention F (M, T, P) holds throughout.

\subsection*{v4.10.0-r59 results, part A: the handedness}
\label{subsec:v4100-r59-results-a}

\paragraph{Theoremlet v4.10.0-r59.1 (Hopf gap as zero point; k245, not blind).}
Replace the area spectrum \(\sqrt{j(j+1)}\) of a puncture in both frames by the semiclassical spectrum \(j+\tfrac12\).  Then the counting
equations close: \(e^{-\pi\gamma_M}=\tfrac12\), \(\gamma_M=\ln2/\pi\), \(\gamma_T=\ln2/(2\pi)\), hence \(\gamma_L=2\) and \(\beta_{\rm sym}=1/\sqrt3\) exactly,
the Hopf value of \texttt{rem:bsym-hopf-gap}.  Without the zero point (spectrum \(j\)) the ratio is \(\log_23\).  The uncertainty enters only as a
zero point, not as a factor.  Reading (not checked): the Hopf construction is the semiclassical limit and the gap \(1/\sqrt3\to0.549\) the quantum
correction of the smallest puncture.  With \(j+\tfrac12\) the data are missed (\(\beta^2=\tfrac13\), \(m_W\)); the result was known before the computation.
\emph{Check.} Wolfram exact, Python control (\texttt{k245}).

\paragraph{Theoremlet v4.10.0-r59.2 (residual rotation as a holonomy; k246, not blind).}
For a three-arc disc with vanishing angular momentum:
(a) on the circle model of 12 November (fixed radius, shape-independent masses) the curvature of the mechanical connection vanishes exactly;
(b) radial and tangential waves with phase offset give a residual rotation \(-4\pi ab/(a^2+2)\) per cycle, maximal exactly at the Koide amplitude
\(a=\sqrt2\) (a restatement of the known \(45^\circ\));
(c) with arc masses \(m\propto\ell^p\) the largest shape cycle in which each arc empties once turns the disc by exactly \(-\pi\), for every tested \(p>0\).
Reading \(\gamma_M\) as the coupling of the connection (Ashtekar--Barbero) gives \(\Delta\Phi=\gamma_M\cdot\tfrac12\), the I1 bridge; \(\gamma_M\) itself does not
arise from the geometry (the curvature contains only \(\pi\) and \(\sqrt2\)).  The open choice moves from ``one light crossing'' to ``largest shape
cycle''; the handedness is set.  \emph{Check.} Wolfram exact, numerical control; a wrong sympy integral (\(-1.26\) instead of \(-\pi\)) was found
and corrected on the way (\texttt{k246}).

\paragraph{Theoremlet v4.10.0-r59.3 (light as mediator; k247--k249, not blind).}
(a) A light signal fixes the time direction \(\tau\) in M; under the tachyon map with holomorphic adjoint, T reads the opposite sense,
\(\tau_T=-\tau_M\) (k247, Wolfram exact).
(b) In the three-well model the threshold to rotation is \(\tfrac23\) of the pendulum frequency, a pure geometric number without a Companion
constant (k248).
(c) In the Julia picture \(c=-1+i\gamma\) the cycle lies on the axis for \(\gamma=0\) (area zero); for \(\gamma\neq0\) an attracting 3-cycle captures every
outer wave and its area has the sign of \(\gamma\) (k249).
(d) The photon point P does not meet the mediator profile of r30.7 (conditions A4, A5).
The handedness is derived nowhere; only threshold and modulus would be derivable, with a definite outer field.  The accompanying picture
``time as residual rotation'' is a reading, not a result.

\paragraph{Theoremlet v4.10.0-r59.4 (handedness from the light weight; k250, not blind).}
The weights \(a,b\) of the light state are diagonal (\(\langle\sigma_z\rangle\)); under orientation reversal (\(\omega\mapsto\omega^2\), \(z\mapsto\bar z\)) they are
unchanged while an orbit sense flips, and under the exchange M\(\leftrightarrow\)T (r81.6(b)) \(a-b\) changes sign.  So \(a-b\) has the parity of \(\tau\),
not of \(\varepsilon\).  For the reference choice r14.2 with \(a\neq b\) the ground state \((x,1,1)\) of \(-(aQ+bR)\) is real and positive for all \(a,b>0\);
the Bargmann invariant \((x^2+2\omega)^3\) has the same phase sign for every real \(x\) (\(\operatorname{Im}=\sqrt3\)), fixed by \(\omega\) alone, and \(K=\tfrac23\)
holds only at \(a=b\).  At the photon point \(\rho=I/2\) no unitary dynamics disturbs \(a=b\).  Balance over r59.1--r59.4: the handedness is set
everywhere; the open choice has arrived at \(\varepsilon\), the choice of \(i\), which over \(\mathbb C\) is a convention.  \emph{Check.} Wolfram symbolic (\texttt{k250}).

\paragraph{Theoremlet v4.10.0-r59.5 (\(\Lambda\) and the rotation rate; k251, not blind).}
(a) Spatial vorticity cannot carry \(\Lambda\): Gödel has \(\Lambda=-\omega^2\) (wrong sign), \(2\omega^2\) dilutes as \(a^{-4}\), and the CMB bound
\((\omega/H)_0<4.7\cdot10^{-11}\) allows at most \(10^{-21}\) of \(\Lambda\).  The rotational energy of an internal field is not \(w=-1\).
(b) In Euclidean de Sitter time is an angle with period \(2\pi/\Omega\) and \(\Lambda=3\Omega^2\) (Gibbons--Hawking, r06).  In the co-rotating frame of a
string with eigenfrequency \(\omega_0\), \(\ddot r=(\Omega^2-\omega_0^2)r\), compared with \(\ddot r=(\Lambda/3)r\) in the de Sitter patch, gives the form
correspondence \(\Lambda/3\leftrightarrow\Omega^2-\omega_0^2\).  The orbit sense enters only as \(s^2=1\).
Verdict: the link is conventional; without an independent fixing of \(\Omega\) or of the string length \(\ell\) it renames \(\Lambda\), as the epoch does
in r56.4.  r56 is unchanged.  \emph{Check.} Wolfram exact (\texttt{k251}).

\paragraph{Theoremlet v4.10.0-r59.6 (H125: relative CP sign of leptons and quarks; stage A confirmed, not blind; stage B open).}
Advance note committed before the computation (\texttt{76f7a6c}).  The neutral family \(su(2)\) of r14.4 (under CT-003) satisfies
\([R_{e_7},R_{e_8}]=-2R_{e_{15}}\) on the lepton block \(H_n\) and \(+2R_{e_{15}}\) on the coloured block; the cause is the non-associativity of the sedenion
right multiplication.  Relative to the Lorentz spin (left action of \(\mathbb H\), Bargmann invariant \(\tfrac14+\tfrac i4\)) the lepton doublet is the
complex conjugate one: the Bargmann invariant of the octant projectors is \(\tfrac14-\tfrac i4\) on the leptons against \(\tfrac14+\tfrac i4\) on the quarks, in
all \(24+24\) assignments of full rank.  The relative sign survives \(i\mapsto-i\), which conjugates both sectors at once; it is the first item of
the series that does not set a sign.  Stage B (blind against future data): under the setting S1 (CP arises in both sectors in the same way from
the family structure) \(\operatorname{sign}J_{\rm PMNS}=-\operatorname{sign}J_{\rm CKM}\), i.e.\ \(\sin\delta_{\rm PMNS}<0\); decided by the first analysis
fixing the sign at \(\ge3\sigma\), deadline 31 December 2035.  Prior probability \(\tfrac12\); a confirmation would mainly support S1.
\emph{Check.} Block computation in numpy on the infrastructure of \texttt{search.py} (deviation from the advance note, stated there); Wolfram exact
for the algebraic core \(B=\tfrac14+\tfrac{is}4\) for \([J_x,J_y]=isJ_z\) (\texttt{k252}).

\subsection*{v4.10.0-r59 results, part B: the Hilbert-shift tower}
\label{subsec:v4100-r59-results-b}

\paragraph{Theoremlet v4.10.0-r59.7 (a generation rule for the tower; k253, not blind).}
(a) HS(0): a real oscillation on a line; embedded in HS(1) it is \(\psi=ae^{i\omega t}+be^{-i\omega t}\) with \(a=b\) (two equal counter-circular parts,
area zero), exactly the time-reversal invariant subfamily.
(b) HS(1): all ellipses with \(a\neq b\); the area \(\pi(a^2-b^2)\) is the \(U(1)\) charge (Wolfram exact); the two circular parts are the sym/antisym
split of the compensator and the half axes \(a\pm b\) the periapsis splitting.  The sense \(u=\pm i\) stays set.
(c) HS(2): \(\mathbb H=\mathbb C^2\), states on the Hopf sphere, Berry phase \(\gamma=-\Omega/2\): the first holonomy, \(SU(2)\), chirality from left
against right multiplication.  HS(3): \(\operatorname{Der}(\mathbb O)=g_2\), stabilizer of a unit \(su(3)\).
(d) \(\dim\operatorname{Der}(A_n)=0,0,3,14,14,14\) for \(n=0,\dots,5\) (Python; Wolfram equal up to \(n=4\)), stabilizer \(0,1,8,8,8\); the circular orbit
closes for every \(u\) only up to \(\mathbb O\) (\(\max|L_u^2+1|\approx10^{-16}\), for the sedenions \(0.85\)), and zero divisors appear (Hurwitz).
HS(6) \(=\mathbb C\otimes\mathbb H\otimes\mathbb O\) carries \(U(1)\), \(SU(2)\), \(SU(3)\) once each on a division level (Dixon).
The groups are symmetries of the field value, not gauge fields; hypercharge and couplings are not derived; ``simpler'' holds algebraically,
not dynamically.

\paragraph{Theoremlet v4.10.0-r59.8 (what HS(4) adds; k254, not blind).}
(a) At the sedenions alternativity, the Moufang identity and norm multiplicativity fail; flexibility and power associativity survive up to
\(n=6\).  The kernel of \(L_a\) on a zero divisor is four-dimensional and its \(G_2\) orbit eleven-dimensional (Moreno).
(b) New is a discrete symmetry: exactly six \(G_2\)-compatible automorphisms (\(S_3\), standard representation on the two copies of the 7, sign on
\(e_8\); Wolfram exact).  Each reflection fixes an octonion subalgebra; the three copies meet only in \(\mathbb R\) and are permuted by \(S_3\).
(c) The largest set of mutually anticommuting unit imaginaries is \(1,3,7,8,9\) for \(n=1,\dots,5\), the Radon--Hurwitz numbers \(\rho(2^n)-1\);
continued \(0,1,3,7\,|\,8,9,11,15\,|\,16,\dots\) (Bott periodicity): HS\((n+4)\) has eight axes more than HS\((n)\).  The tower becomes periodic with
period four; HS(6) has \(11=1+3+7\) axes.  Literature algebra, recomputed.

\paragraph{Theoremlet v4.10.0-r59.9 (r13 reread: the \(S_3\) cycle as a generation step; k255, not blind).}
The r13 result (``the cycle acts under the charge only as a phase'') holds in the eigenbasis of the charge axis \(e_8\) and stands.  In the basis
of the three octonion copies the same \(120^\circ\) rotation is a permutation: on \(V_w=\operatorname{span}\{1,w\otimes f_1,w\otimes f_2\}\) \(S_3\) carries the
permutation representation \(1\oplus2\); the states \(g_k=1/\sqrt3+\sqrt{2/3}\,w\otimes f_k\) are orthonormal, lie in one copy each at the magic angle to
the unit, and are permuted cyclically; in the eigenbasis \(\psi=\operatorname{diag}(1,\omega,\bar\omega)\), the r13 phase.  For
\(\Psi=\alpha+\beta\,w\otimes n(\delta)\) the Koide ratio is \(K=(\alpha^2+\beta^2)/(3\alpha^2)\) (Wolfram exact), so \(K=\tfrac23\) exactly for \(\beta=\alpha\), i.e.\
\(\Psi=\exp(u\pi/4)\) (eighth turn, \(\rho^2=\tfrac12\) as in r12, phase \(\pi/4\) as in H123); then the Brannen form \(1+\sqrt2\cos(\delta-2\pi k/3)\) follows,
with \(\delta\) free.  Limit: the copy charges do not commute and average to \(\tfrac12\mathbb 1\); three generations do not fit into \(\mathbb C\otimes\mathbb S\).
The sedenions supply a generation label and a mixing geometry, the particle content only with three orthogonal carriers.

\paragraph{Theoremlet v4.10.0-r59.10 (mixing from the copy geometry; k256, not blind).}
With the \(S_3\) of the copies as residual symmetry every mixing pattern is fixed before computation: both sectors \(\mathbb Z_3\) give no mixing;
neutrinos \(\mathbb Z_3\) give trimaximal mixing (as H39a, refuted); neutrinos with the copy reflection give \(\theta_{13}=0\) (excluded).  Dictionary
(flavour = copy basis): the TM2 column \((1,1,1)/\sqrt3\) is the real unit, the TM1 column \((2,-1,-1)/\sqrt6\) the imaginary copy axis; exactly
conjugation composed with the copy reflection keeps the TM1 column and returns the CP reading H8.  Nothing beyond known patterns (H4, H8, H39).

\paragraph{Theoremlet v4.10.0-r59.11 (H126: PMNS from the coupling of two octonion directions; not testable).}
Advance note committed before the computation (\texttt{0a1bd3aa}).  Electron and neutrino directions at angle \(\chi\), neutrino state in Koide form,
mass matrix \(M_{jk}=\langle g_j,\Psi_\nu(\Psi_\nu g_k)\rangle\) with \(e_8\) as \(i\).  Structure check (Wolfram exact): \(MM^\dagger\) has eigenvalues
\(\{0,2\cos^2\chi,2\cos^2\chi\}\) for all \(\chi,\delta_\nu\), degenerate everywhere, and the non-associativity contributes exactly zero
(\(M-M_{\rm assoc}\equiv0\)).  By the criterion fixed in advance H126 is \emph{not testable}; the data comparison does not count.  The pairing reads only
\(1\) and \(e_8\), so the associator drops out.  A new attempt needs another pairing, a new number and a new advance note.  (\texttt{k257})

\paragraph{Theoremlet v4.10.0-r59.12 (the tower read backwards; k270--k272, not blind).}
(a) Three colours: the tower offers \(N_c=0,1,3,3,4\) ladders per level (Radon--Hurwitz with one unit fixed).  In the tower model (charge \(k/N_c\),
even sectors left as in r43) only \(N_c=2\) and \(3\) are anomaly free with singlet and fundamental representations; the tower reaches only
\(N_c=3\), first at \(\mathbb O\) (sympy and Wolfram).  HS(2) alone is anomalous (\(\operatorname{Tr}Y^3=\tfrac34\), one doublet, Witten); the tower must
reach colour.
(b) Handedness: on \(\mathbb C\) a genuine bit; from \(\mathbb H\) on it can be turned away by an inner rotation (convention).  What remains a bit is
left/right, which r43 selects by anomaly, up to the frame mirror.
(c) Varying the origin field \(\ddot\psi=g\omega^2\psi\) over all 27 sign patterns of the three shifts: only \((-,-,-)\) gives compact \(G_2/SU(3)\)
without zero divisors; every \(+1\) gives split octonions (leaks from that level, \(sl(3,\mathbb R)\) or \(su(2,1)\) instead of \(SU(3)\)), every \(0\) a
degenerate algebra.  Amplitude, frequency and phase do not act above.  Set remain the r43 frame and M against T; masses, Koide \(\delta\) and the
number of generations are open.

\paragraph{Theoremlet v4.10.0-r59.13 (three generations backwards; k273, not blind).}
A CP phase exists only from \(N=3\) on (\(N=2\): \(J\equiv0\), Wolfram exact); at \(N=3\) every quartet invariant is \(J\cdot\varepsilon\varepsilon\), and \(J\) is a
Bargmann quadrilateral of flavour and mass projectors.  Three is the smallest number with flavour holonomy.  The tower has a cross product only
on \(\operatorname{Im}\mathbb H\) and \(\operatorname{Im}\mathbb O\), a three-dimensional \(\varepsilon\) only on \(\operatorname{Im}\mathbb H\).  Reading: the
generation plane is a further HS(2) level acting \emph{as a vector} (automorphisms \(SO(3)\)).  The doubling units \(e_8,e_{16},e_{32}\) above colour
commute with colour and charge (universality) but act as spin \(\tfrac12\), i.e.\ as a doublet: the known exchange of spin \(\tfrac12\) and spin 1, seen
from the backward side.  Open: a tower level with adjoint HS(2) action (candidate: triality, not computed).

\paragraph{Theoremlet v4.10.0-r59.14 (a nonlinear origin field; k274, not blind).}
With \(\ddot\psi=-\omega^2\psi-\lambda|\psi|^2\psi\) a scale appears (Duffing period \(T\approx2\pi(1-3\lambda A^2/8)\), Wolfram exact via \(K(m)\)); the structure
stays: the term is \(G_2\)-equivariant, the circular orbits \(\exp(u\Omega t)\psi_0\) with \(\Omega^2=\omega^2+\lambda|\psi_0|^2\) are exact up to HS(3)
(residual \(\le5\cdot10^{-13}\)) and break from HS(4) on.  At HS(1) every genuine ellipse precesses (Bertrand); only the line and the circle stay
closed, and the precession follows the orbit sense, so \(\lambda\) selects no handedness.  For \(\lambda<0\) the restoring force changes sign at
\(|\psi|=\omega/\sqrt{|\lambda|}\); beyond it the field lies in the split case of r59.12(c) (reading as a dynamical M/T boundary, not computed).

\paragraph{Remark v4.10.0-r59.15 (balance of the settings).}
Over r59.1--r59.5 the handedness was set at every step; the open choice moved from the light crossing (r58) to the largest shape cycle (r59.2),
to the time direction of the readout (r59.3), and finally to \(\varepsilon\), the choice of \(i\), which is a convention (r59.4).  What is not a
convention is a relative sign of two sectors (r59.6, under CT-003 and the setting S1).  The tower (r59.7--r59.14) makes most choices above HS(0)
forced or conventional; what stays set is the r43 frame, the distinction of M and T, the number of generations and every scale.  No item
derives \(\beta\), a mass, a mixing angle or \(\Lambda\).

\subsection*{v4.10.0-r59 frames}
\label{subsec:v4100-r59-frames}

Part A works in M (light signal, puncture counting) and compares with T only through the tachyon map with holomorphic adjoint (r59.3(a)) or
through \(\gamma_L\) (r59.1); time reversal is complex conjugation, not \(z\mapsto iz\) (r51, r81).  Part B works on the field value in M; the
split case of r59.12(c) and r59.14 is named as a possible reading of T and not computed as such.

\subsection*{v4.10.0-r59 register entries}
\label{subsec:v4100-r59-register}

\begin{center}
\small
\begin{tabular}{@{}>{\raggedright\arraybackslash}p{0.40\textwidth}>{\raggedright\arraybackslash}p{0.52\textwidth}@{}}
\toprule
statement & status after v4.10.0-r59\\
\midrule
r59.1--r59.5, r59.7--r59.10, r59.12--r59.14 & explorations, not blind; recorded\\
H125 stage A (relative orientation of the family \(su(2)\)) & confirmed (not blind)\\
H125 stage B (\(\sin\delta_{\rm PMNS}<0\) under S1) & open, blind against future data, deadline 2035\\
H126 (PMNS from two octonion directions) & not testable (construction degenerate)\\
handedness (absolute) & a convention (choice of \(i\))\\
OP1, OP3, OP4, OP5, Q1a & open\\
\bottomrule
\end{tabular}
\end{center}

\subsection*{v4.10.0-r59 notes}
\label{subsec:v4100-r59-notes}

{\raggedright
\texttt{notes/\allowbreak research/\allowbreak hopf\_\allowbreak gap\_\allowbreak 2026\_\allowbreak 11\_\allowbreak 23/} (k245),
\texttt{holonomie\_\allowbreak restdrehung\_\allowbreak 2026\_\allowbreak 11\_\allowbreak 24/} (k246),
\texttt{licht\_\allowbreak vermittler\_\allowbreak zeit\_\allowbreak 2026\_\allowbreak 11\_\allowbreak 24/} (k247--k249),
\texttt{drehsinn\_\allowbreak lichtgewicht\_\allowbreak 2026\_\allowbreak 11\_\allowbreak 24/} (k250),
\texttt{lambda\_\allowbreak drehrate\_\allowbreak 2026\_\allowbreak 11\_\allowbreak 24/} (k251),
\texttt{h125\_\allowbreak relatives\_\allowbreak cp\_\allowbreak 2026\_\allowbreak 11\_\allowbreak 24/} (k252) with
\texttt{predictions/\allowbreak Vorab\_\allowbreak Notiz\_\allowbreak H125\_\allowbreak Relatives\_\allowbreak CP\_\allowbreak Vorzeichen.md} and the result note,
\texttt{hilbert\_\allowbreak shift\_\allowbreak felder\_\allowbreak 2026\_\allowbreak 11\_\allowbreak 24/} with subfolders \texttt{hs4/} (k254), \texttt{r13/} (k255),
\texttt{mischung/} (k256), \texttt{h126/} (k257, with \texttt{predictions/\allowbreak Vorab\_\allowbreak Notiz\_\allowbreak H126\_\allowbreak Kopplung\_\allowbreak zweier\_\allowbreak Richtungen.md}),
\texttt{rueckwaerts/} (k270--k274).  Each folder holds scripts and their actual outputs.
\par}

\subsection*{v4.10.0-r59 ledger and guard update}
\label{subsec:v4100-r59-ledger-guard}

\[
R_{\rm chk}^{v4100\text{-}r59}=(1,1,0)
\]
for (status recorded; computations checked in the cited scripts; no new computation in this layer; handedness, generations and scales not derived).

The layer spends no guard:
\[
R_{\rm glob}^{(v4100\text{-}r59,+)}=(0,1,1,1,1),
\qquad
\Gamma_{\rm guard}^{v4100\text{-}r59}=(1,1,1,1,1),
\qquad
\beta_4=1/256\ \text{unspent}.
\]
The layer does not imply exact finite \(C5\) zero closure, exact all-readout terminality, \(C6/T2\), a global photon-clock/RH theorem,
or global all-readout closure.  OP1--OP5, C5 and C6/T2 remain open.
